\section{Proof of the main theorem}
\label{sec:conclusion}

The construction uses the same $66\times 66$ pattern $H$ at every depth.
The two preceding estimates now give the required counterexamples.

\begin{proof}[Proof of Theorem~\ref{thm:main}]
For each integer $h\geq1$, set
\[
 \epsilon_h=d_h^{-2}=(386h+2)^{-2},
 \qquad n_h=d_h2^h.
\]
Proposition~\ref{prop:distance} shows that every $H$-free binary matrix
differs from $A_h$ in at least $m^2=\epsilon_h n_h^2$ entries.
Proposition~\ref{prop:copy-count} gives
\[
 \frac{N_H(A_h)}{n_h^{132}}
 \leq \epsilon_h2^{-h}.
\]
These statements prove the asserted construction sequence.

Now fix arbitrary real numbers $c,C>0$.  As $h$ tends to infinity,
\[
 \frac{N_H(A_h)}{\epsilon_h^C n_h^{132}}
 \leq (386h+2)^{2C-2}2^{-h}
 \longrightarrow 0.
\]
Choose an integer $h$ for which the right-hand side is less than $c$.
Then $0<\epsilon_h<1$, the matrix $A_h$ is $\epsilon_h$-far from being
$H$-free, and
$N_H(A_h)<c\epsilon_h^C n_h^{132}$, as required.
\end{proof}

\begin{corollary}[Canonical row--column sampling]
\label{cor:canonical-testing}
Fix $h\ge1$ and an integer $0\le q\le n_h$. Consider a tester that
chooses a uniformly random $q$-element row set $R\subseteq[n_h]$ and,
independently, a uniformly random $q$-element column set
$C\subseteq[n_h]$. It rejects only if the induced submatrix $A_h[R,C]$,
with the inherited orders, contains an ordered copy of $H$. Then
\[
 \Pr(\text{reject on }A_h)\le\epsilon_h2^{-h}q^2,
\]
and the rejection probability is zero when $q<66$. For every fixed
$\rho\in(0,1]$, rejection probability at least $\rho$ requires
\[
 q\ge\sqrt{\rho}\,\epsilon_h^{-1/2}2^{h/2},
\]
hence $q\ge\exp(\Omega(\epsilon_h^{-1/2}))$ along the sequence
$\epsilon_h$ as $h\to\infty$. Here $q$ is the sample size on each axis;
the induced submatrix has $q^2$ entries.
\end{corollary}

\begin{proof}
By Proposition~\ref{prop:copy-count}, every copy of $H$ in $A_h$ uses a
cell of $\mathcal L_h$, and $|\mathcal L_h|=2^h$. For each fixed cell
$(r,c)$, independence gives
$\Pr((r,c)\in R\times C)=(q/n_h)^2$. Rejection implies that
$R\times C$ meets $\mathcal L_h$, so the union bound gives
\[
 \Pr(\text{reject on }A_h)
 \le 2^h(q/n_h)^2=\epsilon_h2^{-h}q^2.
\]
If $q<66$, the sampled submatrix cannot contain $H$. Rearranging the
probability bound gives the necessary sample size, and
$\epsilon_h^{-1/2}=386h+2$ gives its exponential growth along the
stated sequence.
\end{proof}
